\section*{Capítulo \ref{ch:b3:homogeneous-catalysis} --- Catálise homogênea na indústria}

\begin{solution}{exo:b3:homogeneous-catalysis:1}
$\ce{RhCl(PPh3)3}$: 16, Rh(I) $d^8$. $\ce{RhCl(PPh3)2}$: 14, Rh(I). $\ce{RhClH2(PPh3)2}$: 16, Rh(III) $d^6$.
Complexo do alqueno: 18, Rh(III). Alquil-hidreto: 16, Rh(III). A eliminação redutiva
retorna a 14, Rh(I).
\end{solution}

\begin{solution}{exo:b3:homogeneous-catalysis:2}
Butanal, $\ce{CH3CH2CH2CHO}$ (linear), e 2-metilpropanal, $\ce{(CH3)2CHCHO}$ (ramificado).
\end{solution}

\begin{solution}{exo:b3:homogeneous-catalysis:3}
Ciclopent-3-eno-1,1-dicarboxilato de dietila, um anel de cinco membros, com perda de
eteno: os dois grupos $\ce{CH2}$ terminais saem como $\ce{C2H4}$ e os dois carbonos CH internos
se unem.
\end{solution}

\begin{solution}{exo:b3:homogeneous-catalysis:4}
(\textit{E})-3-Fenilprop-2-enoato de metila (cinamato de metila): a arila se adiciona ao
carbono terminal, e a eliminação de $\beta$-hidreto, syn, a partir do confôrmero mais estável,
dá o alqueno \textit{E}.
\end{solution}

\begin{solution}{exo:b3:homogeneous-catalysis:5}
A adição oxidativa é lenta e $\ce{[Rh(CO)2I2]-}$ é o estado de repouso, logo
$v = k[\mathrm{Rh}]_{\mathrm{tot}}[\ce{CH3I}]$. O metanol apenas regenera o iodeto de metila por uma reação rápida com HI;
o iodeto total carregado fixa $[\ce{CH3I}]$, de modo que a velocidade não depende da concentração de
metanol até que reste metanol de menos para converter o HI de volta em
iodeto de metila.
\end{solution}

\begin{solution}{exo:b3:homogeneous-catalysis:6}
Ácido fenilborônico com 4-bromotolueno, ou ácido 4-metilfenilborônico com
bromobenzeno. Ambas funcionam; a primeira usa o ácido borônico mais barato e um brometo de arila
estável e comum.
\end{solution}

\begin{solution}{exo:b3:homogeneous-catalysis:7}
Simetria $C_2$: isotático (os dois sítios selecionam a mesma face). Simetria $C_s$:
sindiotático (sítios especulares se alternam). $\ce{Cp2ZrCl2}$ sem ponte: atático (sítios
aquirais).
\end{solution}

\begin{solution}{exo:b3:homogeneous-catalysis:8}
$\mathrm{TON} = 0.98/0.0010 = 980$; valor médio $\mathrm{TOF} = 980/\qty{2.0}{h} = \qty{490}{h^{-1}}$.
\end{solution}

\begin{solution}{exo:b3:homogeneous-catalysis:9}
A unidade de repetição é $\ce{-[CH=CH-C5H8]-}$, com os dois grupos CH nos carbonos 1 e 3 de um
anel ciclopentano (poli(1,3-ciclopentilenovinileno)). A liberação da
tensão do anel bicíclico torna $\Delta_rH^\circ$ fortemente negativo, o que supera a
perda de entropia translacional; não se forma gás, de modo que a força motriz é
entálpica.
\end{solution}

\begin{solution}{exo:b3:homogeneous-catalysis:10}
$\bar n_n = 1/(1 - p) = 2000$; $M_w/M_n = 1 + p = 1.9995 \approx 2.0$. Um \omterm{def:b3:homogeneous-catalysis:ziegler-natta}{catalisador de Ziegler--Natta} tem
vários tipos de sítios, cada um com seu próprio $p$; a soma de várias distribuições de Schulz--Flory
com médias diferentes é mais larga (dispersidade muitas vezes de 4 a 8).
\end{solution}

\begin{solution}{exo:b3:homogeneous-catalysis:11}
A baixa pressão, $K_1K_2[\ce{H2}] \ll [\mathrm L] + K_1$: $v \approx kK_1K_2[\mathrm{Rh}]_{\mathrm{tot}}[\text{alqueno}][\ce{H2}]/([\mathrm L] + K_1)$,
de primeira ordem em $\ce{H2}$. A alta pressão, o termo $K_1K_2[\ce{H2}]$ domina: $v \to k[\mathrm{Rh}]_{\mathrm{tot}}[\text{alqueno}]$,
independente de $\ce{H2}$. $[\mathrm L]$ só aparece no denominador: a fosfina adicionada diminui $v$
ao empurrar o ródio de volta para o pré-catalisador saturado.
\end{solution}

\begin{solution}{exo:b3:homogeneous-catalysis:12}
A inserção que põe o ródio no carbono interno forma uma alquila secundária
junto a um metal congestionado; as fosfinas volumosas (e um excesso delas, que mantém duas
no metal) tornam esse estado de transição menos favorável que o que dá a alquila
primária, que leva ao aldeído linear. O butanal linear é o que se converte, por
condensação aldólica e hidrogenação, no álcool $C_8$ ramificado usado em
plastificantes; o 2-metilpropanal é um subproduto de menor valor.
\end{solution}

\begin{solution}{pb:b3:homogeneous-catalysis:1}
\textbf{1.} Ir(I), $d^8$: $9 + 4 + 2 + 1 = 16$ elétrons (contagem neutra, com a carga).

\textbf{2.} $\ce{[CH3Ir(CO)2I3]-}$, 18 elétrons, Ir(III).

\textbf{3.} $\ce{[CH3Ir(CO)3I2]}$, 18 elétrons, Ir(III).

\textbf{4.} $\ce{[CH3COIr(CO)2I2]}$, 16 elétrons, Ir(III).

\textbf{5.} $\ce{[CH3COIr(CO)2I3]-}$ (18) elimina $\ce{CH3COI}$, regenerando $\ce{[Ir(CO)2I2]-}$.

\textbf{6.} $\ce{CH3COI + H2O -> CH3COOH + HI}$; $\ce{CH3OH + HI -> CH3I + H2O}$.

\textbf{7.} $\ce{CH3OH + CO -> CH3COOH}$.

\textbf{8.} Com irídio, a inserção migratória é lenta; a adição oxidativa do
iodeto de metila, determinante da velocidade com ródio, é muito mais rápida no irídio, mais
rico em elétrons.

\textbf{9.} Eles retiram um iodeto de $\ce{[CH3Ir(CO)2I3]-}$, dando o tricarbonil neutro, no qual
o metal retrodoa menos e a metila migra mais depressa para um CO mais
eletrofílico.

\textbf{10.} Uma dependência inversa da concentração de iodeto livre: o iodeto desloca
o pré-equilíbrio de volta para o complexo aniônico lento.

\textbf{11.} O metanol só entra depois da etapa determinante da velocidade, convertido em
iodeto de metila por uma reação rápida com HI.

\textbf{12.} $M(\ce{CH3COOH}) = \qty{60.1}{g/mol}$: $n = \qty{5.0e11}{g}/\qty{60.1}{g/mol} = \qty{8.3e9}{mol}$ por ano.

\textbf{13.} $\qty{8.32e9}{mol}/\qty{8000}{h} = \qty{1.04e6}{mol/h}$.

\textbf{14.} $\qty{2.0e5}{L} \times \qty{5.0e-3}{mol/L} = \qty{1.0e3}{mol}$ de irídio, $\qty{1.0e3}{mol} \times \qty{192.2}{g/mol} \approx \qty{190}{kg}$.

\textbf{15.} $\mathrm{TOF} = \qty{1.04e6}{mol/h}/\qty{1.0e3}{mol} = \qty{1.0e3}{h^{-1}}$, cerca de \qty{0.29}{s^{-1}}.

\textbf{16.} $\mathrm{TON} = \qty{8.3e9}{mol}/\qty{1.0e3}{mol} = \num{8.3e6}$ por ano.

\textbf{17.} $\qty{5.0e8}{kg}/(\qty{8000}{h} \times \qty{200}{m^3}) \approx \qty{310}{kg.m^{-3}.h^{-1}}$.

\textbf{18.} $\ce{CO + H2O -> CO2 + H2}$.

\textbf{19.} $1/0.94 = \qty{1.064}{mol}$ de CO por mol de ácido acético.

\textbf{20.} $1.064 - 1 = \qty{0.064}{mol}$ de \ce{CO2} por mol de ácido acético.

\textbf{21.} Uma tonelada corresponde a $\qty{1.0e6}{g}/\qty{60.1}{g/mol} = \qty{1.66e4}{mol}$; $\qty{1.66e4}{mol} \times 0.0638 \times \qty{44.0}{g/mol} \approx \qty{47}{kg}$ de \ce{CO2}.

\textbf{22.} A água é reagente do deslocamento: menos água o torna mais lento e economiza a energia
da secagem do ácido. A água continua necessária para hidrolisar o iodeto de acetila e para manter o
catalisador ativo e dissolvido; com ródio, pouca água precipita o iodeto de
ródio.

\textbf{23.} O hidrogênio: ele se acumula no gás, diminui a pressão parcial de CO
(de modo que é preciso purgar gás, perdendo CO) e pode hidrogenar intermediários em
subprodutos como o ácido propiônico.

\textbf{24.} O catalisador de irídio realiza cerca de $\qty{1.0e3}{h^{-1}}$ ciclos: cada átomo de irídio produz
mil moléculas de ácido acético por hora, cerca de oito milhões por ano.
\end{solution}
